% Part: normal-modal-logic
% Chapter: tableaux
% Section: countermodels

\documentclass[../../../include/open-logic-section]{subfiles}

\begin{document}

\olfileid{nml}{tab}{cou}
      
\olsection{النماذج المضادة من \usetoken{P}{tableau}}

\begin{explain}
  لا يبيّن برهان مبرهنة التمام أنه إذا كان $\Entails !A$، فإن
  $\Proves !A$ فحسب، بل يعطينا أيضًا طريقة لبناء نماذج مضادة لـ~$!A$
  إذا كان $\Entails/ !A$. وفي حالة~$\Log{K}$، تمثل هذه الطريقة
  \emph{إجراء قرار}. افترض، مثلًا، أن $\Entails/ !A$. عندئذ يعطي
  برهان \olref[cpl]{prop:complete-tableau} طريقة لبناء !!{tableau}
  تام، وهذه الطريقة تنتهي دائمًا في الواقع. فقواعد~$\Log{K}$ القضوية
  لا تضيف إلا !!{formula}s ذات بادئات أقل تعقيدًا؛ أي لا يلزم تطبيق
  أي قاعدة قضوية على الفرع أكثر من مرة واحدة لكل صيغة موقعة
  $\sFmla{S}{!A}[\sigma]$. ولا تولّد البادئات الجديدة إلا قاعدة
  \iftag{prvBox}{$\TRule{\False}{\Box}$}{}\iftag{notprvBox,notprvDiamond}{}{
    و}\iftag{prvDiamond}{$\TRule{\True}{\Diamond}$}{}
  \iftag{notprvBox,notprvDiamond}{، ولا يلزم تطبيقها}{، ولا يلزم تطبيقهما}
  إلا مرة واحدة أيضًا (وتنتج بادئة جديدة واحدة).
  وقد يلزم تطبيق
  \iftag{prvBox}{$\TRule{\True}{\Box}$}{}\iftag{notprvBox,notprvDiamond}{}{
    و}\iftag{prvDiamond}{$\TRule{\False}{\Diamond}$}{}
  \iftag{notprvBox,notprvDiamond}{عدة مرات}{عدة مرات}، ولكن مرة واحدة
  فقط لكل بادئة، ولا يُولَّد إلا عدد منته من البادئات الجديدة. ولذلك
  ينتهي البناء بعد عدد منته من المراحل إما إلى فرع مغلق وإما إلى فرع تام.

  متى بُني جدول دلالي له فرع مفتوح تام، أعطانا برهان
  \olref[cpl]{thm:tableau-completeness} نموذجًا صريحًا يرضي المجموعة
  الأصلية من ال!!{formula}s ذات البادئات. ولذلك، ليس الأمر مجرد أنه إذا
  كان $\Gamma \Entails !A$، وُجد !!{tableau} مغلق وكان
  $\Gamma \Proves !A$؛ بل إننا، إذا بحثنا عن ال!!{tableau} المغلق
  بالطريقة المناسبة وانتهينا إلى !!{tableau} «تام»، فلن نعلم فقط أن
  $\Gamma \Entails/ !A$، بل سنتمكن فعلًا من بناء نموذج مضاد.
\end{explain}

\iftag{prvBox}{
\begin{ex}
  نعلم أن $\Proves/ \Box(p \lor q) \lif (\Box p \lor \Box q)$.
  يبدأ بناء جدول دلالي هكذا:
  \begin{oltableau}
    [\pFmla{\False}{\Box(p \lor q) \lif (\Box p \lor \Box q)}{1},
      just = \TAss, checked
      [\pFmla{\True}{\Box(p \lor q)}{1},
        just = {\TRule{\False}{\lif}[1]}, 
        [\pFmla{\False}{\Box p \lor \Box q}{1},
          just = {\TRule{\False}{\lif}[1]}, checked
          [\pFmla{\False}{\Box p}{1},
            just = {\TRule{\False}{\lor}[3]}, checked
            [\pFmla{\False}{\Box q}{1},
              just = {\TRule{\False}{\lor}[3]}, checked
              [\pFmla{\False}{p}{1.1}, 
                just = {\TRule{\False}{\Box}[4]}, checked
                [\pFmla{\False}{q}{1.2}, 
                  just = {\TRule{\False}{\Box}[5]}, checked
                ]
              ]
            ]
          ]
        ]
      ]
    ]
  \end{oltableau}
  لم يكتمل ال!!{tableau} بعدُ بالطبع. وفي الخطوة التالية ننظر في السطر
  الوحيد الذي لا يحمل علامة فحص: ال!!{formula} ذات البادئة
  $\sFmla{\True}{\Box(p \lor q)}[1]$ في السطر~$2$. ويقتضي بناء
  الجدول الدلالي التام تطبيق قاعدة $\TRule{\True}{\Box}$ على كل بادئة
  مستعملة على الفرع، أي على كل من $1.1$ و$1.2$:
  \begin{oltableau}
    [\pFmla{\False}{\Box(p \lor q) \lif (\Box p \lor \Box q)}{1},
      just = \TAss, checked
      [\pFmla{\True}{\Box(p \lor q)}{1},
        just = {\TRule{\False}{\lif}[1]}, 
        [\pFmla{\False}{\Box p \lor \Box q}{1},
          just = {\TRule{\False}{\lif}[1]}, checked
          [\pFmla{\False}{\Box p}{1},
            just = {\TRule{\False}{\lor}[3]}, checked
            [\pFmla{\False}{\Box q}{1},
              just = {\TRule{\False}{\lor}[3]}, checked
              [\pFmla{\False}{p}{1.1}, 
                just = {\TRule{\False}{\Box}[4]}, checked
                [\pFmla{\False}{q}{1.2}, 
                  just = {\TRule{\False}{\Box}[5]}, checked
                  [\pFmla{\True}{p \lor q}{1.1}, 
                    just = {\TRule{\True}{\Box}[2]}
                    [\pFmla{\True}{p \lor q}{1.2}, 
                      just = {\TRule{\True}{\Box}[2]}
                    ]
                  ]
                ]
              ]
            ]
          ]
        ]
      ]
    ]
  \end{oltableau}
  والآن لا تحمل الأسطر 2 و8 و9 علامات فحص. لكن لم تُضف أي بادئة
  جديدة، فنطبق $\TRule{\True}{\lor}$ على السطرين~8 و~9 في جميع الفروع
  الناتجة (ما دامت لا تنغلق):
  \begin{oltableau}
    [\pFmla{\False}{\Box(p \lor q) \lif (\Box p \lor \Box q)}{1},
      just = \TAss, checked
      [\pFmla{\True}{\Box(p \lor q)}{1},
        just = {\TRule{\False}{\lif}[1]}, checked
        [\pFmla{\False}{\Box p \lor \Box q}{1},
          just = {\TRule{\False}{\lif}[1]}, checked
          [\pFmla{\False}{\Box p}{1},
            just = {\TRule{\False}{\lor}[3]}, checked
            [\pFmla{\False}{\Box q}{1},
              just = {\TRule{\False}{\lor}[3]}, checked
              [\pFmla{\False}{p}{1.1}, 
                just = {\TRule{\False}{\Box}[4]}, checked
                [\pFmla{\False}{q}{1.2}, 
                  just = {\TRule{\False}{\Box}[5]}, checked
                  [\pFmla{\True}{p \lor q}{1.1}, 
                    just = {\TRule{\True}{\Box}[2]}, checked
                    [\pFmla{\True}{p \lor q}{1.2}, 
                      just = {\TRule{\True}{\Box}[2]}, checked
                      [\pFmla{\True}{p}{1.1},
                        just = {\TRule{\True}{\lor}[8]}, checked, close
                      ]
                      [\pFmla{\True}{q}{1.1},
                        just = {\TRule{\True}{\lor}[8]}, checked
                        [\pFmla{\True}{p}{1.2},
                          just = {\TRule{\True}{\lor}[9]}, checked]
                        [\pFmla{\True}{q}{1.2},
                          just = {\TRule{\True}{\lor}[9]}, checked, close]
                      ]
                    ]
                  ]
                ]
              ]
            ]
          ]
        ]
      ]
    ]
  \end{oltableau}
  يبقى فرع مفتوح واحد، وهو تام. ومنه نعرّف نموذجًا عوالمه
  $W = \{1, 1.1, 1.2\}$ (وهي البادئات الوحيدة الواقعة على الفرع
  المفتوح)، وعلاقة الإتاحة فيه
  $R = \{\tuple{1, 1.1}, \tuple{1, 1.2}\}$، وإسناده $V(p) =
  \{1.2\}$ (لأن السطر~11 يحتوي $\sFmla{\True}{p}[1.2]$) و
  $V(q) = \{1.1\}$ (لأن السطر~10 يحتوي
  $\sFmla{\True}{q}[1.1]$). ويُرسم النموذج في
  \olref{fig:counter-Box}، ويمكنك التحقق من أنه نموذج مضاد لـ
  $\Box(p \lor q) \lif (\Box p \lor \Box q)$.
  \begin{figure}
  \begin{center}
    \begin{tikzpicture}[modal]
      \node[world] (w1) [label={[align=right]right:\mFalse{p}\\ \mFalse{q}}]{$1$}; 
      \node[world] (w2) [label={[align=right]right:\mFalse{p}\\ \mTrue{q}},
        above left=of w1]{$1.1$}; 
      \node[world] (w3) [label={[align=right]right:\mTrue{p}\\ \mFalse{q}},
        above right=of w1] {$1.2$};
      \draw[->] (w1) to (w2);
      \draw[->] (w1) to (w3);
    \end{tikzpicture}
  \end{center}
  \caption{نموذج مضاد لـ$\Box(p \lor q) \lif (\Box p \lor \Box q)$.}
  \ollabel{fig:counter-Box}
\end{figure}
\end{ex}          
}
{
\begin{ex}
  نعلم أن $\Proves/ (\Diamond p \land \Diamond q) \lif \Diamond(p
  \land q)$. يبدأ بناء جدول دلالي هكذا:
  \begin{oltableau}
    [\pFmla{\False}{(\Diamond p \land \Diamond q) \lif \Diamond(p \land q)}{1},
      just = \TAss, checked
      [\pFmla{\True}{\Diamond p \land \Diamond q}{1},
        just = {\TRule{\False}{\lif}[1]}, checked
        [\pFmla{\False}{\Diamond(p \land q)}{1},
          just = {\TRule{\False}{\lif}[1]}
          [\pFmla{\True}{\Diamond p}{1},
            just = {\TRule{\True}{\land}[2]}, checked
            [\pFmla{\True}{\Diamond q}{1},
              just = {\TRule{\True}{\land}[2]}, checked
              [\pFmla{\True}{p}{1.1}, 
                just = {\TRule{\True}{\Diamond}[4]}, checked
                [\pFmla{\True}{q}{1.2}, 
                  just = {\TRule{\True}{\Diamond}[5]}, checked
                ]
              ]
            ]
          ]
        ]
      ]
    ]
  \end{oltableau}
  لم يكتمل ال!!{tableau} بعدُ بالطبع. وفي الخطوة التالية ننظر في السطر
  الوحيد الذي لا يحمل علامة فحص: ال!!{formula} ذات البادئة
  $\sFmla{\False}{\Diamond(p \land q)}[1]$ في السطر~$3$. ويقتضي بناء
  الجدول الدلالي التام تطبيق قاعدة $\TRule{\False}{\Diamond}$ على كل
  بادئة مستعملة على الفرع، أي على كل من $1.1$ و$1.2$:
  \begin{oltableau}
    [\pFmla{\False}{(\Diamond p \land \Diamond q) \lif \Diamond(p \land q)}{1},
      just = \TAss, checked
      [\pFmla{\True}{\Diamond p \land \Diamond q}{1},
        just = {\TRule{\False}{\lif}[1]}, checked
        [\pFmla{\False}{\Diamond (p \land q)}{1},
          just = {\TRule{\False}{\lif}[1]}
          [\pFmla{\True}{\Diamond p}{1},
            just = {\TRule{\True}{\land}[2]}, checked
            [\pFmla{\True}{\Diamond q}{1},
              just = {\TRule{\True}{\land}[2]}, checked
              [\pFmla{\True}{p}{1.1}, 
                just = {\TRule{\True}{\Diamond}[4]}, checked
                [\pFmla{\True}{q}{1.2}, 
                  just = {\TRule{\True}{\Diamond}[5]}, checked
                  [\pFmla{\False}{p \land q}{1.1}, 
                    just = {\TRule{\False}{\Diamond}[3]}
                    [\pFmla{\False}{p \land q}{1.2}, 
                      just = {\TRule{\False}{\Diamond}[3]}
                    ]
                  ]
                ]
              ]
            ]
          ]
        ]
      ]
    ]
  \end{oltableau}
  والآن لا تحمل الأسطر 3 و8 و9 علامات فحص. لكن لم تُضف أي بادئة
  جديدة، فنطبق $\TRule{\False}{\land}$ على السطرين~8 و~9 في جميع
  الفروع الناتجة (ما دامت لا تنغلق):
  \begin{oltableau}
    [\pFmla{\False}{(\Diamond p \land \Diamond q) \lif \Diamond(p \land q)}{1},
      just = \TAss, checked
      [\pFmla{\True}{\Diamond p \land \Diamond q}{1},
        just = {\TRule{\False}{\lif}[1]}, checked
        [\pFmla{\False}{\Diamond(p \land q)}{1},
          just = {\TRule{\False}{\lif}[1]}
          [\pFmla{\True}{\Diamond p}{1},
            just = {\TRule{\True}{\land}[2]}, checked
            [\pFmla{\True}{\Diamond q}{1},
              just = {\TRule{\True}{\land}[2]}, checked
              [\pFmla{\True}{p}{1.1}, 
                just = {\TRule{\True}{\Diamond}[4]}, checked
                [\pFmla{\True}{q}{1.2}, 
                  just = {\TRule{\True}{\Diamond}[5]}, checked
                  [\pFmla{\False}{p \land q}{1.1}, 
                    just = {\TRule{\False}{\Diamond}[3]}, checked
                    [\pFmla{\False}{p \land q}{1.2}, 
                      just = {\TRule{\False}{\Diamond}[3]}, checked
                      [\pFmla{\False}{p}{1.1},
                        just = {\TRule{\False}{\land}[8]}, checked, close
                      ]
                      [\pFmla{\False}{q}{1.1},
                        just = {\TRule{\False}{\land}[8]}, checked
                        [\pFmla{\False}{p}{1.2},
                          just = {\TRule{\False}{\land}[9]}, checked]
                        [\pFmla{\False}{q}{1.2},
                          just = {\TRule{\False}{\land}[9]}, checked, close]
                      ]
                    ]
                  ]
                ]
              ]
            ]
          ]
        ]
      ]
    ]
  \end{oltableau}
  يبقى فرع مفتوح واحد، وهو تام. ومنه نعرّف نموذجًا عوالمه
  $W = \{1, 1.1, 1.2\}$ (وهي البادئات الوحيدة الواقعة على الفرع
  المفتوح)، وعلاقة الإتاحة فيه
  $R = \{\tuple{1, 1.1}, \tuple{1, 1.2}\}$، وإسناده $V(p) =
  \{1.1\}$ (لأن السطر~6 يحتوي $\sFmla{\True}{p}[1.1]$) و
  $V(q) = \{1.2\}$ (لأن السطر~7 يحتوي
  $\sFmla{\True}{q}[1.2]$). ويُرسم النموذج في
  \olref{fig:counter-Diamond}، ويمكنك التحقق من أنه نموذج مضاد لـ
  $(\Diamond p \land \Diamond q) \lif \Diamond (p \land q)$.
  \begin{figure}
  \begin{center}
    \begin{tikzpicture}[modal]
      \node[world] (w1) [label={[align=right]right:\mFalse{p}\\ \mFalse{q}}]{$1$}; 
      \node[world] (w2) [label={[align=right]right:\mTrue{p}\\ \mFalse{q}},
        above left=of w1]{$1.1$}; 
      \node[world] (w3) [label={[align=right]right:\mFalse{p}\\ \mTrue{q}},
        above right=of w1] {$1.2$};
      \draw[->] (w1) to (w2);
      \draw[->] (w1) to (w3);
    \end{tikzpicture}
  \end{center}
  \caption{نموذج مضاد لـ$(\Diamond p \land \Diamond q) \lif
    \Diamond (p \land q)$.}
  \ollabel{fig:counter-Diamond}
\end{figure}
\end{ex}          
}

\end{document}
